% Generated draft: requires mathematical and source review.
\documentclass[11pt]{article}
\usepackage[margin=25mm]{geometry}
\usepackage{fontspec}
\setmainfont{Latin Modern Roman}
\usepackage{amsmath,amssymb,mathtools,unicode-math}
\setmathfont{Latin Modern Math}
\usepackage{enumitem,needspace,xurl,hyperref}
\setlength{\parindent}{0pt}
\setlength{\parskip}{5pt}
\setlength{\emergencystretch}{3em}
\newcommand{\N}{\mathbb{N}}
\newcommand{\Z}{\mathbb{Z}}
\newcommand{\Q}{\mathbb{Q}}
\newcommand{\R}{\mathbb{R}}
\newcommand{\C}{\mathbb{C}}
\newcommand{\F}{\mathbb{F}}
\newcommand{\A}{\mathbb{A}}
\newcommand{\PP}{\mathbb{P}}
\newcommand{\E}{\mathbb{E}}
\newcommand{\eps}{\varepsilon}
\newcommand{\dd}{\,\mathrm{d}}
\begin{document}
\section*{UnsolvedMath ID: 30006153}
\textbf{Source problem code:} catalogue entry
\subsection{Definitions and mathematical statement}
The supplied record attributes both assertions below to Robertson (1994), as Conjectures 5 and 6. Its dated curation describes the projective-plane assertion as solved and the general assertion as open; these status claims have not been independently verified here.

For a provisional conventional formulation, let graphs be finite undirected graphs, allowing loops and parallel edges, and let surfaces be compact connected surfaces without boundary, including both orientable and nonorientable surfaces. These conventions are not specified in the supplied extract and require confirmation against the source. No cellularity assumption on embeddings is stated in the extract.

A graph is \emph{$2$-edge-connected} if it is connected and deleting any one edge leaves it connected. For an embedding \(\iota:G\hookrightarrow S\), define its \emph{face-width} by
\[
\operatorname{fw}(\iota)=\min_C\bigl|C\cap\iota(G)\bigr|,
\]
where \(C\) ranges over noncontractible simple closed curves in \(S\) whose intersection with the embedded graph is finite. Curves may pass through vertices; each intersection point is counted once. If there is no noncontractible simple closed curve, set \(\operatorname{fw}(\iota)=+\infty\).

A \emph{nowhere-zero $4$-flow} on \(G\) consists of an orientation of its edges and a function
\[
f:E(G)\longrightarrow\{-3,-2,-1,1,2,3\}
\]
satisfying conservation at every vertex:
\[
\sum_{e\in\delta^+(v)}f(e)=\sum_{e\in\delta^-(v)}f(e)
\qquad(v\in V(G)).
\]
Here \(\delta^+(v)\) and \(\delta^-(v)\) denote outgoing and incoming edges, respectively; an oriented loop contributes to both sums.

\textbf{Projective-plane assertion (source Conjecture 5).} For every $2$-edge-connected graph \(G\) and every embedding \(\iota:G\hookrightarrow\mathbb{RP}^{2}\),
\[
\operatorname{fw}(\iota)\ge4
\quad\Longrightarrow\quad
G\text{ has a nowhere-zero }4\text{-flow}.
\]

\textbf{General surface assertion (source Conjecture 6).} For every surface \(S\), there exists a positive integer \(r(S)\) such that, for every $2$-edge-connected graph \(G\) and every embedding \(\iota:G\hookrightarrow S\),
\[
\operatorname{fw}(\iota)\ge r(S)
\quad\Longrightarrow\quad
G\text{ has a nowhere-zero }4\text{-flow}.
\]
The threshold depends only on \(S\), not on \(G\) or its embedding. Taking a positive integer threshold expresses the source's unspecified constant threshold equivalently because finite face-width is integer-valued. The graph and surface conventions above, and whether embeddings must be cellular, remain source-verification issues rather than settled hypotheses of the supplied extract.
\subsection{Short English statement}
Robertson's general assertion asks whether each surface has a face-width threshold forcing every $2$-edge-connected graph embedded in that surface to admit a nowhere-zero $4$-flow. The record also includes the projective-plane assertion with threshold $4$. Supplied curation dated 2026-08-22 labels that special case solved and the general assertion open; neither status is independently verified here. The intended graph, surface, and embedding conventions require confirmation.
\subsection{Sources}
\begin{itemize}
\item Local UnsolvedMath registry: supplied complete record for permanent ID \(30006153\), exact source problem code \texttt{catalogue entry}, titled ``Robertson's Surface Four-Flow Conjecture,'' in the Oberwolfach Reports Open Problems collection.
\item Supplied report citation: \emph{Graph Theory} (2025), Oberwolfach Reports, DOI 10.4171/owr/2025/1, pp.\ 5--74. \url{https://doi.org/10.4171/owr/2025/1}. The extract attributes Conjectures 5 and 6 to Robertson (1994); the record's proposed-year field is 2025, corresponding to the report year rather than the attributed conjecture year.
\item Supplied literature source, labeled ``projective-plane special case'': \url{https://arxiv.org/abs/2405.16586}. No authors, title, or theorem locator for this reference are supplied.
\item Literature assessment checked at \texttt{2026-08-22}: the supplied curation labels the projective-plane case solved and the all-surfaces assertion unresolved. This is supplied dated curation, not an independently verified literature assessment. The background explicitly requests independent verification of time-sensitive status and cited claims.
\item The registry labels the overall problem \texttt{open} and its statement \texttt{corrected\_verified}. Its verification note says that only formula delimiters, notation, grammar, and typography were repaired while preserving mathematical content and scope; this does not establish mathematical correctness.
\end{itemize}
\end{document}
